\section{Discussion and limitations}
\label{sec:discussion}

\paragraph{What the speech buys.}
On LongVideoBench, the confirmed gain over the starting design is a package effect that the post hoc decomposition traces mostly to the change of frames, and the tuned pipeline showed no demonstrated advantage over R3, which uses no transcript; that comparison changes more than the transcript, so it does not isolate a transcript effect in either direction.
The lecture results look different, although we ran no formal test of the difference between the two datasets.
The S2m package reads the whole transcript and raises the share of questions with a candidate inside the labelled source-slide interval from 81.2\% to 89.0\%, which is where most of its hit@4 difference from S1 arises.
Yet R5, a training-free rule that spreads subtitle similarity over time and fuses it with option-aware visual similarity, reaches the labelled slide at a similar rate (79.8\% against 79.4\%), and the two are shown neither to differ nor to be equal.
Removing Head A's learned residual gave no confirmed change on LongVideoBench, and training on lecture questions gave no confirmed change against training on CG-Bench.
A plausible reading, which our design does not isolate, is that at this budget much of the benefit comes from searching the whole transcript and video with the answer options in the query, and that a small learned correction adds little beyond what careful rules provide, at least for the rules and models we tested.

\paragraph{Reading accuracy with evidence measures.}
For the two declared lecture comparisons, accuracy and the evidence proxy lead to the same reading (\cref{sec:lectures}), but accuracy alone cannot show why a system fails.
The answer model already gets 51.2\% of the lecture questions right with no video at all, and generated lecture questions can often be answered from the transcript.
The slide intervals show where the labelled slide is missed, though they do not establish why either.
With S2m's candidate pools held fixed, a perfect picker could raise hit@4 by at most 9.6 points.
No comparable bound follows for a different candidate stage, which would change both the pool and what the picker does with it, and the 9.6-point figure does not bound accuracy.
Building on earlier work that evaluates evidence and not only answers~\cite{li2025kfsbench,huang2026egvqa}, we suggest that frame-selection studies report, next to accuracy, an evidence-placement measure, no-video controls, the realized frame and text budget of every system, and confidence intervals that resample whole videos.

\paragraph{Limitations.}
Our lecture labels are automatic slide intervals with known false matches, and the questions, for training as well as for evaluation, are generated by a model from slides.
A hit means that a shown frame's timestamp falls in a labelled interval, not that the slide or the needed fact was visible.
The human check covers 100 sampled questions and has the weaknesses listed in \cref{sec:human}: one reviewer accepted every question, 33 questions had a single reviewer, independence rests on the organizer's account, and no reviewer judged the frames the systems chose.
Several contrasts compare packages that change more than one component.
Search work is not matched either: on the lecture test, R3--R5 decode up to 96 ranked times, while S1 and the heads decode at most about 54.
R5 is our adaptation of VSI, tuned over ten settings on development data, so stronger published methods remain untested under our conditions.
The lecture study uses a single 4B answer model and M3AV's shipped transcripts rather than a local speech recognizer, so we do not show robustness to recognition errors.
The decomposition, the residual ablation and the loss analysis were run after the relevant test set was revealed and are post hoc.

Finally, we record three deviations from our own plans, all documented in the supplement.
First, each answer model's ledger had its own 20,000-call cap, so the declared total of 20,000 calls for the LongVideoBench confirmation was not enforced; the ledgers charged 25,373 calls, including API follow-ups on the earlier locked test.
Second, one crashed control answer in the lecture run was retried against our rule; scored as the rule requires, every full-cohort accuracy is unchanged and the transcript-only-wrong subset has 322 rather than the as-run 323 questions (\cref{sec:lectures}).
Third, in the frozen lecture answer run the observed reserved memory exceeded the card's capacity, and one out-of-memory crash and restart is recorded (\cref{sec:cost}).
None of these changes a primary result, though each remains a deviation.

\section{Conclusion}
\label{sec:conclusion}
With at most four frames and a frozen small answer model, a transcript-guided pipeline improves accuracy over our starting design as a package.
In lectures, a pipeline built around a speech-to-moment locator places frames inside the labelled source-slide interval more often than a tuned rule-based pipeline.
It shows no demonstrated advantage over a training-free rule that fuses subtitle and visual scores, and two small trained heads add no confirmed accuracy on fresh long-video questions.
We think the more lasting lesson concerns measurement.
Accuracy alone could not have told us at which stage the labelled slide was missed, or how often the model answered correctly with no video at all; evidence-placement measures and no-video controls could, and we hope future studies of frame selection will report them.
